\section{Introduction}\label{sec:introduction}

A Riemannian metric is Einstein if its Ricci tensor is a constant
multiple of the metric. This fixes the trace of the sectional curvatures
while allowing individual curvatures to vary. The classification
conjecture in dimension four asks whether the round $S^4$, the round
$\mathbb{RP}^4$, and Fubini--Study $\mathbb{CP}^2$ exhaust the connected
closed Einstein four-manifolds with strictly positive sectional curvature,
up to positive scaling and isometry. Yang explicitly lists these three
models in Conjecture~1 of~\cite{Yang2000}. We prove the classification
in the closed setting.

\begin{theorem}\label{thm:classification}
Let $(M,g)$ be a connected smooth closed Einstein four-manifold with
strictly positive sectional curvature, without an orientability
assumption. There is a constant $a>0$ such that $(M,ag)$ is isometric to
the standard round $S^4$, to $\mathbb{CP}^2$ with its Fubini--Study
metric, or to $\mathbb{RP}^4$ with its standard round metric.
\end{theorem}

Strict sectional positivity makes the Ricci tensor positive definite,
so the Einstein constant is positive. The oriented classification is the
principal task. Until the final subsection of
Section~\ref{sec:classification}, we therefore work on a connected
oriented smooth closed four-manifold, choose an orientation, and rescale
the metric so that $\operatorname{Ric}=3g$. The final orientation-cover
argument uses the oriented conclusion proved there and determines the
possible nonorientable quotient.

\paragraph{Earlier results.}
The problem combines rigidity of the Einstein equation with the global
restrictions of positive sectional curvature. Berger's early rigidity
work under strict quarter pinching~\cite{Berger1961} was followed by
quantitative sectional-curvature bounds of Yang, Costa, and
Cao--Tran~\cite{Yang2000,Costa2004,CaoTran2018}. These bounds impose
additional quantitative hypotheses, some expressed as upper bounds
relative to the Einstein constant rather than as pinching ratios.
Gursky--LeBrun prove the Fubini--Study conclusion for compact oriented
Einstein manifolds with nonnegative sectional curvature and a nonzero
positive-definite intersection form~\cite[Theorem~A]{GurskyLeBrun1999}.
Koca proves the Fubini--Study conclusion for a complete positively sectionally
curved Einstein metric Hermitian for an integrable complex
structure~\cite[Theorem~2]{Koca2014}. Under the different condition of
nonnegative isotropic curvature, the four-dimensional local-symmetry
result is due to Micallef--Wang~\cite[Theorem~4.4]{MicallefWang1993};
Brendle extends the Einstein rigidity result to higher dimensions
\cite[Theorem~1]{Brendle2010}.

Recent results impose other geometric or topological conditions.
Dameno obtains half-conformal flatness from fibrewise constancy of the
scalar curvature of an associated twistor metric, and obtains the
standard models in the compact positive-scalar case
\cite[Theorem~1.1 and Corollary~1.2]{Dameno2026}.
Catino--Dameno use separation of the largest and middle eigenvalues
of one Weyl block to obtain conformal K\"ahler geometry after at most a double cover
\cite[Theorems~1.1 and~1.4]{CatinoDameno2026}.
Di Cerbo treats the closed oriented nonflat case with nonnegative
sectional curvature and equality everywhere in a Weyl-norm bound,
obtaining a locally symmetric $S^2\times S^2$ universal cover
\cite[Proposition~3.2]{DiCerbo2026}.

Write $\chi$ for the Euler characteristic and $\tau$ for the signature.
Cheng proves the oriented positive-sectional classification under
$\chi\le3$, and consequently in the presence of a nontrivial Killing
field~\cite[Theorem~1.2 and Corollary~1.4]{Cheng2026}. His argument
combines radial Riccati comparison with Weyl-norm gaps and
characteristic-number identities. Gursky--Malchiodi prove the classification for closed oriented Einstein
four-manifolds with positive sectional curvature under
$2\chi-3|\tau|\le4$, using a Green-function comparison and ADM mass
\cite[Theorem~1.1]{GurskyMalchiodi2026}. These cited results retain
additional curvature, geometric, topological, or symmetry hypotheses.
The argument below addresses the interaction of the two Weyl blocks
without such an additional restriction.

\paragraph{The two-block problem.}
The Hodge star on two-forms has eigenvalues $+1$ and $-1$; its two
rank-three eigenbundles are denoted by $\Lambda^+$ and $\Lambda^-$.
In the normalization $\operatorname{Ric}=3g$, the curvature operator
restricts to $\operatorname{Id}+T$ and $\operatorname{Id}+U$ on these
bundles, where $T,U$ are symmetric trace-free endomorphisms. They are
the two Weyl curvature blocks. The Einstein equation removes the
off-diagonal Ricci block but does not force either Weyl block to vanish.
Our central intermediate conclusion is \emph{half-conformal flatness}:
\[
 T\equiv0\qquad\text{or}\qquad U\equiv0.
\]

Use Hilbert--Schmidt norms and put
\[
 v=|T|/\sqrt6,\qquad b=|U|/\sqrt6,\qquad
 x=-\lambda_{\min}(T),\qquad y=-\lambda_{\min}(U).
\]
The exact sectional-positivity condition is $x+y<2$. The inequalities
$v\le x\le2v$ and $b\le y\le2b$ imply $v+b<2$, but this weaker norm
condition alone is insufficient for the tensor estimate below.
Gursky--LeBrun use the corresponding nonstrict cone and norm inequalities
and analyze the case of two nonzero Weyl blocks
\cite[Lemma~1 and Theorem~B]{GurskyLeBrun1999}.
For a function $f$, write $\mathbf E f$ for its integral divided by total
Riemannian volume. The gap and its parallel equality case are established
in Gursky--LeBrun~\cite[Theorem~1]{GurskyLeBrun1999}: a nonzero block
satisfies $\mathbf E v^2\ge1$, or the exchanged inequality for $b$.
We derive this statement in our normalization, including zeros of a block.

Scalar coupling of the two blocks also has a clear antecedent. Yang
chooses a constant $t$ for which $|W^+|-t|W^-|$ has mean zero and combines
the Weyl equation, refined Kato inequality, and Poincar\'e inequality
\cite[\S6]{Yang2000}. Cao--Tran similarly balance regularized powers of
the norms~\cite[Proposition~2.2]{CaoTran2018}. Here a volume and moment
argument forces equality of the two second moments, and a weighted
Hessian square supplies the additional tensor estimate.

\paragraph{From moment balance to half-conformal flatness.}
The characteristic identities relate the two second moments to volume
and to the dimensions of harmonic self-dual and anti-self-dual two-form
spaces. An upper polar-volume deficit, the two gaps, and a strengthened
scalar variance inequality force
$\mathbf E v^2=\mathbf E b^2$ when both blocks are nonzero
(Proposition~\ref{prop:balanced-moments}). This balance controls the mean
of $v-b$, enabling the variance inequality to contribute the sign needed
in the final integral argument.

The main differential estimate uses the Bianchi identity twice. It first
places the gradient of a Weyl block in one irreducible tensor type.
The curvature commutator then determines all but one component of its
full, generally nonsymmetric Hessian. A curvature-weighted quadratic
form is positive on the remaining component because its lower bound is
$4-2x-2y>0$. Completing its square with different scalar weights for the
two blocks and integrating by parts produces inequalities involving only
first derivatives (Proposition~\ref{prop:weighted-hessian}). The proof
retains both the derivatives of the scalar weights and the derivative of
the curvature operator acting on the matrix slot. This local estimate
does not use moment balance.

The two weighted inequalities, a scalar identity, and the variance bound
produce an integrand $\mathcal I$ with nonnegative mean. Explicit rational
comparison functions give $c\ge1/10$, nonnegative quantities $r,j$ built
from the norm gradients, and a scalar $\Theta$ such that
\[
 \mathbf E\mathcal I\ge0,\qquad
 \mathcal I\le-\frac{(\Theta-cr)^2}{c}-c(j-r^2),
 \qquad r^2\le j.
\]
The pointwise bounds follow from the weighted Hessian calculation and
exact finite polynomial signs. Moreover, $r^2<j$ almost everywhere that
either norm gradient is nonzero. Integration therefore makes both norms
constant. If both blocks were nonzero, their separate gaps would then
contradict $v+b<2$. This proves half-conformal flatness.

The final geometric argument has a different role. Once one block
vanishes, the characteristic identities leave two harmonic counts. The
exceptional count forces $s_0=\mathbf E(v^2+b^2)=3$ and volume
$4\pi^2/3$; a conditional lower polar-volume estimate excludes precisely
that case. Equality in the surviving Weyl gap makes the curvature
parallel. Matching its tensor with a standard model and using simple
connectivity then gives a global isometry. Figure~\ref{fig:proof-map}
records the assembly.

\begin{figure}[htbp]
\centering
\begin{tikzpicture}[
 box/.style={draw=blue!40!black,rounded corners=2pt,fill=blue!3,
   text width=4.15cm,align=center,inner sep=5pt,font=\footnotesize},
 arr/.style={-{Latex[length=2mm]},draw=blue!55!black,thick},
 note/.style={font=\scriptsize,fill=white,inner sep=2pt,align=center}]
 \node[box] (scalar) at (0,0) {Scalar and upper volume estimates\\Sections~\ref{sec:analytic}--\ref{sec:volume}};
 \node[box] (hessian) at (6.2,0) {Weighted Hessian inequality\\Section~\ref{sec:hessian}};
 \node[box] (balance) at (0,-1.55) {Equal second moments\\Section~\ref{sec:volume}};
 \node[box] (integral) at (6.2,-1.55) {$\mathbf E\mathcal I\ge0$ for the two blocks\\Section~\ref{sec:coupling}};
 \node[box] (signs) at (0,-3.1) {Exact pointwise bounds\\Sections~\ref{sec:coupling}--\ref{sec:certificates}};
 \node[box] (vanish) at (6.2,-3.1) {Half-conformal flatness\\Strictness in Section~\ref{sec:classification}};
 \node[box] (lower) at (0,-4.65) {Conditional\\lower volume bound\\Section~\ref{sec:classification}};
 \node[box] (parallel) at (6.2,-4.65) {Parallel curvature and model tensor\\Section~\ref{sec:classification}};
 \node[box] (simply) at (0,-6.2) {Simple connectivity\\Appendix~\ref{app:topology}};
 \node[box] (global) at (6.2,-6.2) {Global model isometry\\Section~\ref{sec:classification}};
 \draw[arr] (scalar)--(balance);
 \draw[arr] (hessian)--(integral);
 \draw[arr] (balance)--node[note,above]{variance\\term}(integral);
 \draw[arr] (integral)--(vanish);
 \draw[arr] (signs)--node[note,above]{pointwise\\signs}(vanish);
 \draw[arr] (vanish)--(parallel);
 \draw[arr] (lower)--(parallel);
 \draw[arr] (parallel)--(global);
 \draw[arr] (simply)--(global);
\end{tikzpicture}
\caption{Assembly of the proof. Moment balance enters the nonnegative
integral; the exact pointwise bounds and active-gradient strictness force
one-block vanishing. The conditional lower bound identifies the remaining
parallel tensor, and simple connectivity permits global continuation.}
\label{fig:proof-map}
\end{figure}


The radial estimates use discrete Jacobi determinants and retain the
order of curvature matrices along each geodesic. The lower refinement is
conditional on the moment value just described. The article gives the
geometric and analytic arguments, together with the finite polynomial
formulas and sign proofs. The accompanying checker reproduces the saved
finite arithmetic within the precise scope stated in its coverage record.

\paragraph{Organization.}
Section~\ref{sec:preliminaries} fixes curvature conventions and derives
the differential identities, including their behavior at zeros.
Section~\ref{sec:analytic} proves the weighted variance and Weyl-gap
estimates and states the characteristic identities.
Section~\ref{sec:volume} proves the upper deficit and moment balance.
Section~\ref{sec:hessian} derives the weighted Hessian inequality, and
Section~\ref{sec:coupling} combines it with the scalar estimates.
Section~\ref{sec:certificates} proves the resulting finite signs.
Section~\ref{sec:classification} proves one-block vanishing, the
conditional lower volume bound, and the full metric classification.
Appendix~\ref{app:analytic} contains the deferred analytic and heat
proofs, while Appendix~\ref{app:topology} proves the required topology.
